% MAIN SNAPSHOT — LSPV/RMV::CP22LIFT1::2026-10-13T01:40:00Z
% !TEX program = xelatex
\documentclass[11pt]{article}
\usepackage[T1]{fontenc}
\usepackage{lmodern,seqsplit,amsmath,amssymb,amsthm,booktabs,tabularx,array,xurl,hyperref}
\usepackage[margin=0.87in]{geometry}
\hypersetup{hidelinks,pdftitle={CP22LIFT1: Cutting Planes Transport through Exact-(2,2) Normalization}}
\setlength{\emergencystretch}{3em}
\newtheorem{theorem}{Theorem}
\newtheorem{lemma}[theorem]{Lemma}
\newtheorem{corollary}[theorem]{Corollary}
\newcommand{\CP}{\mathsf{CP}}
\newcommand{\PC}{\mathsf{PC}}
\newcommand{\TT}{\mathcal T}
\newcommand{\poly}{\operatorname{poly}}
\title{CP22LIFT1\\\large Two-Way Cutting Planes Proof Transport\\through the Actual Selector Exact-$(2,2)$ Compiler}
\author{UNIVERSAL\_CORE\_SYNTHESIS / LSPV-5}
\date{13 October 2026, 01:40 UTC}
\begin{document}
\maketitle
\begin{center}
\fbox{\parbox{0.94\linewidth}{\textbf{P versus NP remains OPEN.}
The proof below establishes a model-specific two-way proof translation for the
actual PC22\_ALLFIELD1 selector compiler. Its status is
\textbf{PROVED/EXACT at the declared Cutting Planes scope, independent review
pending}. This is not a transfer from arbitrary SAT executions. The
experiment checks explicit integer certificates, not a proof-assistant
formalization. The all-field Polynomial Calculus lower bound and the existing
counting shortcut are retained with their original attribution.}}
\end{center}

\section{Result, exact quantifiers, and scope}
A proper 3-clause has three distinct underlying variables. Let $F$ be a
nonempty ordered list of $m$ proper 3-clause occurrences, densely numbered
on its active variables, and let $L=|\operatorname{enc}(F)|$ in binary indexed
clause-list encoding. Repeated source clauses are allowed. Let $\TT(F)$ be
\emph{the selector-and-duplication compiler of PC22\_ALLFIELD1}, not the later
neutral-slot ALLFIELD22 compiler. Its output has
\[
                 v=72m,\qquad M=96m,\qquad
                 |\operatorname{enc}(\TT(F))|=\Theta(m\log(m+2)).       \tag{1}
\]
The output is proper, clause-distinct, signed exact-$(2,2)$; it is not
asserted clause-monotone or minimally unsatisfiable.

Cutting Planes, abbreviated $\CP$, is not Polynomial Calculus, abbreviated
$\PC$. A $\CP$ line is $a\cdot x\ge b$, with integer coefficients and
Boolean variables. Axioms are clause inequalities $\sum_{\ell\in C}\ell\ge1$
and $x\ge0$, $-x\ge-1$, where a negative literal denotes $1-x$.
Rules are sum, nonnegative integer scaling, and
\[
 \frac{a\cdot x\ge b}{(a/d)\cdot x\ge\lceil b/d\rceil}
 \quad(d\ge1,\ d\mid a_i\text{ for every }i).                         \tag{2}
\]
We count the complete binary encoding of every line and inference, including
coefficient bits, variable indices and parent references. This is the standard
integer-inequality model of \cite[Section 2.1]{ln}, with explicit indexing
costs. Proofs use the variables of the formula in question.

\begin{theorem}[\texttt{CP22-TWO-WAY-PROOF-TRANSPORT1}]\label{thm:main}
There are two fixed deterministic transducers $\mathsf{Lift}$ and
$\mathsf{Project}$ and a fixed polynomial $p$ such that, for every such $F$:
\[
\begin{aligned}
 \pi:F\vdash_{\CP}0\ge1
 &\Longrightarrow
 \mathsf{Lift}(F,\pi):\TT(F)\vdash_{\CP}0\ge1,\\
 \tau:\TT(F)\vdash_{\CP}0\ge1
 &\Longrightarrow
 \mathsf{Project}(F,\tau):F\vdash_{\CP}0\ge1.
\end{aligned}                                                        \tag{3}
\]
Time, peak workspace, and explicit output length of each translation are at
most $p(L+|\pi|)$ or $p(L+|\tau|)$, respectively. The compiler map is computed
from the supplied $F$, not acquired by an oracle or inferred for free from an
anonymized target. A supplied signed permutation and clause/literal
reordering can also be transported with polynomial overhead.
\end{theorem}
Thus $F$ and $\TT(F)$ have polynomially related $\CP$ proof sizes. In
particular, a suitable source $\CP$ lower bound can be transported into this
exact-$(2,2)$ syntax. This does not assert that a source lower bound against
all algorithms exists, or that arbitrary SAT deciders emit $\CP$ proofs.

\section{The actual compiler and its recovered dependencies}
For a source clause $(a_j\lor b_j\lor c_j)$, with letters denoting literals,
introduce a common selector $s$ and fresh $y_j,z_j$ and emit
\begin{align*}
 (s\lor a_j\lor y_j),&\quad(\neg y_j\lor b_j\lor c_j),\\
 (\neg s\lor\neg a_j\lor z_j),&\quad
 (\neg z_j\lor\neg b_j\lor\neg c_j).
\end{align*}
Call this list $U(F)$ and duplicate the whole occurrence list to get
$B(F)=U(F)\uplus U(F)$. Intermediate deduplication is forbidden. Every
variable has equally many positive and negative occurrences, at least two of
each. Pair opposite occurrences in list order, with one new clone per pair.
Join each variable's clones in a directed implication cycle.

There are $R=12m$ clones and $R$ implications. Pad each triple of
implications with the three fresh forced-zero ports of one copy of the
recovered gadget $Q$. It has 15 variables and 19 clauses, so (1) follows.
The full fixed gadget appears in the supplied program. Its all-zero model
and three forced-zero ports are the same dependency as in PCRETRACT1.
The existing width-three Resolution port proof has a 69-node ancestor
subproof, including the 19 axioms; our program reconstructs and replays this
finite proof, rather than assuming semantic implications are free.

The previously proved projection identifies clones and sets
\[
            s=0,\quad y_j=\neg a_j,\quad z_j=0,
            \quad\text{all gadget variables}=0.                       \tag{4}
\]
It sends the second split clause to the source clause and all other split,
cycle, and gadget clauses to Boolean-valid clauses or inequalities. All
properness, occurrence, and semantic compiler facts here are adopted from
PC22\_ALLFIELD1, not relabelled as new discoveries \cite{pc22,pcretract}.

\section{Proof of the forward translation}
\subsection{Derive ports and clone equalities}
Resolution of $A\lor x$ with $B\lor\neg x$ is simulated in $\CP$ as follows.
Sum the two clause inequalities and cancel $x+(1-x)=1$. Each literal of the
resolvent $A\cup B$ now has multiplicity one or two. Add the nonnegativity
bound of each multiplicity-one literal to make every multiplicity two, then
divide by two with rounding. This derives the exact resolvent inequality.
For the fixed port certificate, this is a constant amount of work.

Hence each padded implication gives
\[
 (-u+v+p\ge0)+(-p\ge0)\quad\Longrightarrow\quad -u+v\ge0.             \tag{5}
\]
On every clone cycle, fix its first clone $r$ as representative. Prefix sums
along the cycle give $-r+u\ge0$; suffix sums give $-u+r\ge0$.
There are $O(R)$ such sparse lines in total. Adding the appropriate
equalities to each transformed source clause reconstructs its inequality
on representatives. This obtains the first copy of $U(F)$ without
materializing pairwise equality for every pair of clones.

Summing the first two split clauses eliminates $y_j$ and gives
\[
             s+a_j+b_j+c_j\ge1.                                     \tag{6}
\]
Summing the other pair similarly gives
\[
        (1-s)+(1-a_j)+(1-b_j)+(1-c_j)\ge1.                            \tag{7}
\]
The remaining step is a fully charged proof-by-cases simulation. Naively
carrying a selector coefficient through (2) is not sufficient.

\subsection{The guard-rounding lemma}
\begin{lemma}[Explicit guard repair]\label{lem:guard}
Suppose each source axiom $A(x)\ge b$ is available as $A(x)+\ell\ge b$,
where $\ell$ is one fresh Boolean literal, and all source Boolean bounds
remain available. A $\CP$ refutation of the source can be transformed in
polynomial encoded size into a derivation of $\ell\ge1$.
\end{lemma}
\begin{proof}
For every source line $a\cdot x\ge b$, maintain a line
$a\cdot x+k\ell\ge b$ with an explicitly stored nonnegative integer $k$.
A clause axiom has $k=1$ and a Boolean bound has $k=0$. Addition and
nonnegative scaling update $k$ by the same operations.

For a source division by $d$, let
\[
            K=d\lceil k/d\rceil.                                    \tag{8}
\]
Add $(K-k)\ell\ge0$ using the Boolean nonnegativity bound. Now all
coefficients are divisible by $d$, including the guard, and division yields
\[
        (a/d)\cdot x+\lceil k/d\rceil\ell\ge\lceil b/d\rceil.         \tag{9}
\]
For $\ell=1-s$, the normalized right side is $b-K$ before division;
$K$ is divisible by $d$, so
$\lceil(b-K)/d\rceil=\lceil b/d\rceil-K/d$. Thus (9) is valid for
both signs of the guard, not just for $s$.

The final line is $k\ell\ge1$. The case $k=0$ would be an unconditional
contradiction derived from the guarded axioms. Those axioms are satisfiable
when $\ell=1$, since each clause inequality then permits all Boolean source
assignments. Soundness excludes this case. Divide by $k>0$ to get
$\ell\ge1$.

If the source proof has encoded size $S$, every guard coefficient has
$O(S)$ bits: additions increase bit length by at most one, scalings add the
encoded multiplier length, and rounding in (8) adds at most the encoded
divisor length. There are only constantly many added operations per source
inference. All arithmetic and output therefore have polynomial bit cost.
\end{proof}

Apply the lemma to the supplied source proof using (6) and $\ell=s$.
Apply it a second time to the same proof after complementing every source
variable, using (7) and $\ell=1-s$. Signed substitution respects (2), because
any induced constant shift at a divisible line is itself divisible by $d$.
We obtain $s\ge1$ and $-s\ge0$, whose sum is $0\ge1$.
The port and clone preparation has $O(m)$ sparse lines; the guarded
simulations have polynomial bit size in $S$. This proves the first half of
Theorem~\ref{thm:main}.

\section{Projection back, including a repaired failed shortcut}
Apply (4) to every line of a target proof. A literal/constant substitution
maps an integer affine expression to another such expression. Addition and
nonnegative scaling commute with it. At a division line, every original
coefficient is divisible by $d$. Variable identification preserves this
property, and the constant shift $c$ introduced by substitution is divisible
by $d$. Therefore
\[
                   \lceil(b-c)/d\rceil=\lceil b/d\rceil-c/d.          \tag{10}
\]
So the projected inference is legal in the identical integer proof system.

A projected initial axiom is either a source clause inequality or a
Boolean-valid affine inequality. The latter must be derived, not silently
deleted. For $a\cdot x\ge b$, Boolean validity is equivalent to
\[
                     b\le\sum_{a_i<0} a_i.                          \tag{11}
\]
Scaled Boolean bounds derive $a\cdot x\ge\sum_{a_i<0}a_i$.
Adding the required multiple of $0\ge-1$, itself obtained from
$x\ge0$ and $-x\ge-1$, yields the exact weaker right side $b$.
These have polynomial encoded size. The final contradiction remains
$0\ge1$, proving the second half of Theorem~\ref{thm:main}.

\paragraph{Archived failed implementation.}
The first projection prototype treated every non-source image as the zero
inequality. Exact replay rejected this. For example, setting $z=0$ in
$(\neg z\lor\neg b\lor\neg c)$ with positive $b,c$ gives
$-b-c\ge-2$, not $0\ge0$. The final implementation uses (11).
This repairs our prototype; it does not refute the predecessor's
\emph{Polynomial Calculus} theorem, which has a different axiom encoding.
The second rejected shortcut was division when $d\nmid k$; (8) supplies the
missing guard slack explicitly.

\section{Concrete application: all-field algebraic hardness, short integer proofs}
Let $G$ be a simple 4-regular bipartite graph with $h$ vertices on each side,
and let $G'$ add one previously absent edge. The subset-cardinality source
$F_h=\mathrm{SC}(G')$ has one variable per edge, all positive triples at each
left vertex, and all negative triples at each right vertex. Thus
\[
       n_h=4h+1,\quad m_h=8h+12,\quad
       v_h=576h+864,\quad M_h=768h+1152.                             \tag{12}
\]
For the expansion hypotheses of Mik\v{s}a--Nordstr\"om
\cite[Theorem 4.8]{mn}, source PC/PCR degree is $\Omega(h)$ over every
fixed field. The already established PC22\_ALLFIELD1 projection and the
standard size--degree bound give
\[
       \operatorname{Size}_{\PC/\mathrm{PCR}}(\TT(F_h))
       \ge 2^{\Omega(v_h)}=2^{\Omega(L_h/\log L_h)}.                  \tag{13}
\]
This is an inherited theorem with graph hypotheses, not a conclusion from
our finite graphs. In contrast, the following $\CP$ proof is explicit.

For four incident positive edge variables, sum their four triple clauses:
$3\sum x_i\ge4$. Divide by three to obtain $\sum x_i\ge2$.
For five incident variables, derive this bound on each four-element subset,
sum the five bounds and divide by four:
\[
                        4\sum_{i=1}^5 x_i\ge10
                        \quad\Longrightarrow\quad\sum_{i=1}^5 x_i\ge3.
\]
Apply the same reasoning to negative literals on the right. Every right
vertex permits at most two selected edges. Summing left quotas and right
capacities gives
\[
                 \sum_{e\in E(G')}x_e\ge2h+1,
                 \qquad \sum_{e\in E(G')}x_e\le2h,                  \tag{14}
\]
a contradiction. There are $O(h)$ source proof lines, with polynomial
bit size. Theorem~\ref{thm:main} lifts (14) through the entire selector
compiler, giving $O(v_h)$ target lines and a safe
$O(v_h^2\log(v_h+2))$ encoded-size bound.

The counting escape itself was already in PC22\_ALLFIELD1, and the newer
neutral-slot ALLFIELD22 has a stronger source-free bounded-local terminal
confirmed by REDTEAM-STRENGTHEN8 \cite{rt8}. Our contribution is the
explicit two-way $\CP$ proof translation for this exact selector compiler,
not a new general counting algorithm or an unannotated proof-search result.
Lauria--Nordstr\"om provides a related PC-hard/CP-easy separation for a
\emph{different graph-colouring construction} \cite[Proposition 1.3]{ln}.

\section{Reproducible evidence and limits}
The standard-library generator and separately coded exact-integer replay
checker processed every axiom, Boolean bound, sum, scaling and division.
Each main batch also transported the proof through a genuine random signed
variable permutation and independent clause/literal reordering, projected
the full target proof back to the source, rejected a corrupted final line,
and verified a satisfying lift for the graph without the extra edge.
\begin{center}\small
\begin{tabular}{rrrrrr}
\toprule
$h$ & Target vars & Clauses & Canonical lines & Projected lines & Seconds\\
\midrule
5 &3744&4992&62223&62434&1.77\\
6 &4320&5760&71779&72022&2.10\\
8 &5472&7296&90891&91198&2.56\\
12&7776&10368&129115&129550&3.86\\
\bottomrule
\end{tabular}
\end{center}
Scrambled proofs have the same line counts as canonical proofs. These four
batches replayed 1,063,220 target/scrambled/projected lines. Seven additional
controls replayed 291,408 lines: properized Horn and 2SAT sources, odd-charge
$K_4$ Tseitin, all eight proper three-variable clauses, independent products,
duplicate source occurrences, and clause/literal order changes. Total final
replay count is \textbf{1,354,628 lines}, with zero failures. Source proofs
and internal repeated checker calls are not added again to that number.

The even controls verified 23,808 target clauses against explicit satisfying
assignments. Each main batch checked Q's 244 models and 3,825 exact guard
rounding identities; 2,525 tested guard/divisor pairs per batch disallow the
naive unadjusted division. The seven-control batch took 10.82 seconds. Every
invocation used a 30-second external timeout; no batch exceeded it.
Properizing Horn or 2SAT need not preserve those syntactic subclasses, so
these names describe the source origins only. The finite circulant graphs
are not certified asymptotic expanders. No minimum proof size, general SAT
runtime, or proof-assistant theorem was measured.

\paragraph{Evidence files.}
\path{cp22_audit.py}, \path{cp22_controls.py}, \path{h5.json},
\path{h6.json}, \path{h8.json}, \path{h12.json}, \path{controls.json}, and
\path{h5_certificate.json}. The latter contains all 62,223 explicit target
proof lines and their generating inferences. The main script SHA-256 is
\begin{center}\small\ttfamily
\seqsplit{d49869d2faa31bae336df16c1e6aa13b6c64728d93802bf98e3b89c6bfb6ddc9}
\end{center}
Commands and the complete artifact hash manifest are supplied in the bundle.

\section{SAME\_REALIZATION, barriers, and the next proof target}
The source proof, compiler map, target proof generator, and independent
checker are explicit separate procedures. Their parsing, clause writing,
coefficient arithmetic, indexing, equality derivations, guard repair,
projection, and proof output all have charged polynomial bit cost. No part
of their cost is attributed to an arbitrary pre-existing one-bit decider.
The theorem provides no sublinear-space implementation, fixed favorable time
exponent, shrinking successor, or FLMV landing. Both actual manuscript
drivers retain their all-input acquisition and hereditary-access gaps.

This is a propositional proof-system simulation, not an oracle separation.
No nonrelativizing or nonalgebrizing hinge is supplied. A construction using
integer inequalities should not be called an algebrization escape just
because it differs from finite-field Polynomial Calculus. No Natural Proofs
claim is invoked: there is no proposed large constructive circuit property.
The exact cross-domain import is integer optimization and rounding; parity,
Schur-square, cryptographic binding, and black-box extraction analogies do
not prove an ordinary-machine bound.

\paragraph{Quantifier clarification.}
BLACKBOXQ1's equivalence identifies a separation-strength obligation.
Equivalence to $P\ne NP$ does not itself refute a candidate theorem or make
it unusable in a proof. What is invalid is treating that theorem as already
proved by bookkeeping or by a bound in a different model. The direct target
remains a bound for every uniform deterministic decider of a fixed NP
language; no subsequent broadening to general SAT would be necessary.

\paragraph{One next check.}
Use the exact backward $\CP$ projection as a gate on the next proposed source:
seek a proper 3-CNF family with a verified superpolynomial $\CP$ bound as well
as the desired algebraic bound, then check whether its encoding admits this
literal/constant projection. This is a concrete stronger-proof-system
benchmark task, not the universal-computation theorem. The latter remains
\textbf{OPEN} and must be tracked separately.

\section{Corpus, integration, and publication record}
\paragraph{CORPUS REVIEWED.}
Read current AGENTS, CODEX and repository skill; current root synthesis
opening; actual main4/main5 driver openings and the current ALLFIELD22 audit
include; BLACKBOXQ1; the four specialist reports named by that snapshot
(TRACESIMQ1, LOGFPRINT1, EARLY22RED1, PCRETRACT1); the substantive
PC22\_ALLFIELD1 and REDTEAM-STRENGTHEN8 proofs; and the EARLY22 governance
memo. Queried repository content for prior Cutting Planes retraction. This
is a scoped review, not a claim to have exhausted the repository or all
current specialist updates. No local checkout was available, so repository
preflight and whole-root compilation were not run.

\paragraph{DUPLICATION / CONTRADICTION AUDIT.}
The actual general selector normalizer, all-field PC/PCR source theorem,
counting shortcut, Q port theorem, and later neutral-slot orbit terminal are
prior work. They are adopted, not re-proved as novel discoveries. The new
proof-transport claim does not contradict any of them. The initial local
projection bug and illegal guard divisions are archived above. No historical
minimality or continuation claim is silently strengthened or retired.

\paragraph{EARLY22-INTEGRATION1 acknowledgment.}
\textbf{ALREADY INTEGRATED, with a new model-specific extension.} This run
read the governance memo and existing BLACKBOXQ1 integration ledger, not the
original Drive excerpts again. Preserve separate entries for the compiled
exact-$(2,2)$ Resolution-hard MU family, the genuinely clause-monotone
cubic-pair theorem/finite calibrations, and the affine-continuation
limitation. The latter retains the ledger's supersession by 21 P C B; the
older EARLY22RED1 request to reopen the other two deletion types is not
adopted. Exact new manuscript location: this report, Theorem~\ref{thm:main}
and Sections 6--8. Existing four-specialist acknowledgment statuses are not
upgraded without new acknowledgments. One next check is the independent
review of guard rounding and all target axiom images in the projection.

\paragraph{WHAT CHANGED RELATIVE TO ROOT MAIN.}
Add the two-way $\CP$ proof-transport entry and the explicit short-$\CP$
certificate for the selector-compiled subset-cardinality family. Preserve
BLACKBOXQ1 and the three EARLY22 entries. Main4 remains ARC and main5 remains
WRAPPER. An owner-ready insertion and peer handoff accompany this report.
Publication, readback, and root-sync status are recorded separately; this
scientific body does not assert an unperformed root update.

\begin{thebibliography}{9}
\bibitem{pc22} LSPV-4. PC22\_ALLFIELD1: Universal Exact-$(2,2)$ Normalization
and Field-Independent Algebraic Bounds. Canonical repository report,
13 October 2026.
\path{research/agents/lspv-4/2026-10-13_0010_LSPV-4_PC22_ALLFIELD1.tex}.
\bibitem{pcretract} LSPV-4. PCRETRACT1: Algebraic Degree Retraction of the
Exact-$(2,2)$ Compiler. Canonical repository report, 13 October 2026.
\path{research/agents/lspv-4/2026-10-13_0849_LSPV-4_PCRETRACT1.tex}.
\bibitem{mn} M. Mik\v{s}a and J. Nordstr\"om.
A Generalized Method for Proving Polynomial Calculus Degree Lower Bounds.
CCC 2015, LIPIcs 33, 467--487. Full-length arXiv version 1505.01358,
Theorems 2.2 and 4.8, Definition 4.6.
\url{https://arxiv.org/abs/1505.01358}.
\bibitem{ln} M. Lauria and J. Nordstr\"om.
Graph Colouring is Hard for Algorithms Based on Hilbert's Nullstellensatz
and Gr\"obner Bases. arXiv:2306.00125. Section 2.1 and Proposition 1.3.
\url{https://arxiv.org/abs/2306.00125}.
\bibitem{rt8} REDTEAM-STRENGTHEN8. ALLFIELD22 Survives Its Local Audit but
Does Not Supply a WRAPPER Successor. 13 October 2026.
\path{research/agents/proof/2026-10-13_0113_REDTEAM-STRENGTHEN8.tex}.
\bibitem{blackbox} LSPV-5. BLACKBOXQ1. Canonical integration and
quantifier ledger, 13 October 2026.
\path{research/main/2026-10-13_0948_MAIN_BLACKBOXQ1.tex}.
\end{thebibliography}
\end{document}
% END_OF_VERIFIED_MAIN::CP22LIFT1
